% !TeX program = pdflatex
% ALIUS Bulletin standalone facsimile source.
% Compile from the ALIUS-bulletin project root for exact visual fidelity.
% If the shared PDF asset is absent, this file falls back to extracted text.
\providecommand{\ALIUSRootPrefix}{}

\ifdefined\ALIUSIssueBuild
  \IfFileExists{\ALIUSRootPrefix Shared-assets/original-pdfs/Issue06/alius-bulletin-issue-06-doss-cinco-sins-of-psychedelic-research.pdf}{%
    \includepdf[pages=-,fitpaper=true]{\ALIUSRootPrefix Shared-assets/original-pdfs/Issue06/alius-bulletin-issue-06-doss-cinco-sins-of-psychedelic-research.pdf}%
  }{%
    \typeout{ALIUS warning: missing source PDF Shared-assets/original-pdfs/Issue06/alius-bulletin-issue-06-doss-cinco-sins-of-psychedelic-research.pdf}%
  }%
  \expandafter\endinput
\fi

\documentclass[a4paper,11pt]{article}
\usepackage[margin=22mm]{geometry}
\usepackage{pdfpages}
\usepackage[hidelinks]{hyperref}
\usepackage{microtype}
\usepackage{parskip}
\title{On the Cinco Sins of Psychedelic Research}
\author{An interview with Manoj Doss by George Fejer}
\date{ALIUS Bulletin 6}
\begin{document}
\IfFileExists{\ALIUSRootPrefix Shared-assets/original-pdfs/Issue06/alius-bulletin-issue-06-doss-cinco-sins-of-psychedelic-research.pdf}{%
  \includepdf[pages=-,fitpaper=true]{\ALIUSRootPrefix Shared-assets/original-pdfs/Issue06/alius-bulletin-issue-06-doss-cinco-sins-of-psychedelic-research.pdf}%
}{%
  \maketitle
  \section*{Metadata}
  \textbf{Piece type:} commentary\par
  \textbf{DOI:} 10.5281/zenodo.10797875\par
  \textbf{Keywords:} Cognitive Neuropsychopharmacology, Psychedelic Research, Bias, Subjectivity\par
  \section*{Abstract}
  In this interview, Manoj Doss addresses the "Cinco Sins of Psychedelic Research," highlighting methodological concerns and biases, including those stemming from self-experimentation and researchers' personal drug experiences. Advocating for a closer integration of cognitive psychology and neuropharmacology, he calls for empirical approaches to study psychedelics' effects on cognition and memory. Doss's critique underscores the importance of behavioral measures to advance a scientific understanding of psychedelic substances.\par
  \section*{Extracted Text}
  \begingroup
  \small
  \noindent On The Cinco Sins of Psychedelic Research An interview with Manoj Doss By George Fejer Manoj Doss manoj.doss@austin.utexas.edu Department of Psychiatry \& Behavioral Sciences, University of Texas, Austin, USA George Fejer george.fejer@uni-konstanz.de Department of Cognitive Psychology University of Constance, Konstanz, Germany Cite as: Doss, M., Fejer, G., (2024). On The Cinco Sins of Psychedelic Research. An interview with Manoj Doss by George Fejer, ALIUS Bulletin, 6, https://doi.org/https://doi.org/10.5281/zenodo.10797875 Abstract In this interview, Manoj Doss addresses the "Cinco Sins of Psychedelic Research," highlighting methodological concerns and biases, including those stemming from self-experimentation and researchers' personal drug experiences. Advocating for a closer integration of cognitive psychology and neuropharmacology, he calls for empirical approaches to study psychedelics' effects on cognition and memory. Doss's critique underscores the importance of behavioral measures to advance a scientific understanding of psychedelic substances. Keywords: Cognitive Neuropsychopharmacology, Psychedelic Research, Bias, Subjectivity Hello Manoj, thank you for taking the time to do this interview. I wanted to use this opportunity to ask you some questions, given your reputation for having a critical voice within the psychedelic research community. You have raised some sharp concerns over issues such as personal biases and methodological rigor, as outlined in your lectures on the "Cinco Sins of Psychedelics" (e.g., Morris, 2022). In your line of research, you have investigated the effects of psychedelics and various psychoactive drugs on episodic memory, among many other interesting topics. But I wanted to open up the conversation by discussing your job title, which you have ascribed as being a "Cognitive Neuropsychopharmacologist". What does this entail? Basically, I have been using that title as a joke to give myself the longest career title possible. Some researchers call themselves Cognitive Neuroscientists, others call themselves Neuropsychopharmacologists, but I ALIUS Bulletin n6 (2022) aliusresearch.org/bulletin Doss - On The Cinco Sins of Psychedelic Research identify as both since I am a cognitive psychologist who measures the brain while under the influence of different drugs. Can you give an example of the type of research questions that motivate you in trying to bring these two disciplines closer together? Are you more interested in characterizing the effects of drugs by looking at how they affect memory, or in probing the underlying processes of memory through drug interventions? What started your initial interest in entering the field? Most pharmacologists tend to stay clear of cognitive psychology, yet I believe this field holds significant value, particularly when examining drugs with bizarre mental effects by investigating how they affect information processing. What started my initial interest was the investigation of episodic memory, a subdomain of cognitive psychology that has been studied for over a century. It is possibly one of the oldest research topics in the field, and as such, it has developed stringent methods. But cognitive psychologists typically avoid psychopharmacology, because drug interventions are considered as a very messy type of whole-brain manipulation. However, in my opinion, this is an excellent place to examine pitfalls in the field of psychopharmacology, as I have discovered that many researchers are not well-trained in measuring cognitive processes like memory. It can be difficult to be good at everything, but I believe that cognitive psychology offers a lot of valuable tools. On the flip side, the field of cognitive psychology can often be pedantic, with researchers debating the smallest effect sizes which do not necessarily have any immediate implications for improving something in the world. Pursuing a career in psychopharmacology next to cognitive psychology was a good way for me to broaden my research interests and make them more relevant. One of the other major benefits of cognitive psychology [...] is that you can try to set yourself up to fail. [...] In contrast, psychedelic drug studies ... investigating subjective effects do not take the risk of producing null results. ALIUS Bulletin n6 (2024) aliusresearch.org/bulletin Doss - On The Cinco Sins of Psychedelic Research To date ,the most high-impact paper that I have published was also the simplest data set to analyze and collect. It demonstrated that when people retrieve memories under the influence of THC, they are more likely to have a false memory due to an increased false alarm rate (Doss et al., 2018). There was no fancy analysis or brain imaging in this paper, but it was accepted for publication in the journal Biological Psychiatry because no one in the field had ever properly isolated the drug effects in relation to the retrieval phase of a memory test. I believe there is a lot of such low-hanging fruit, where simple behavioral paradigms can go a long way in explaining how drugs affect behavior and cognition. That may sound opportunistic, but I am in fact very interested in differentiating the effects of different drugs using precise cognitive manipulations. [Y]ou cannot really use drugs a lot of times to show how the mind works. I was trying to prove him wrong at first, but now I am beginning to believe he is correct. One of the other major benefits of cognitive psychology, in my opinion, is that you can try to set yourself up to fail. A researcher can run an experiment without necessarily expecting to find any effect, but if they find an effect nonetheless, that's probably a good sign. In contrast, many psychedelic drug studies set themselves up for success because they are all about investigating subjective effects and do not take the risk of producing null results. People report a wide range of subjective effects, so the researchers create a questionnaire to assess things like mystical experiences, administer the drug alongside the questionnaire, and discover that people score these higher while under the influence of the drug. What else would they anticipate was happening? Did they expect to find any negative results? If you give people a lot of different questionnaires while under the influence of drugs, they will probably score higher on at least some of them. To me, it would seem much more interesting to read a study that produced unexpected results, and using methods of cognitive psychology is a good way to do that. ALIUS Bulletin n6 (2024) aliusresearch.org/bulletin Doss - On The Cinco Sins of Psychedelic Research Is there something particularly insightful about using drug interventions in the field of cognitive psychology, or more specifically, memory research? Charan Ranganath, my advisor at UC Davis, once told me that you cannot really use drugs a lot of times to show how the mind works (although he does not remember saying that). I was trying to prove him wrong at first, but now I am beginning to believe he is correct. I have not learned much about the mind from studying drugs, but I am still trying. There are numerous confirmatory findings of things we already knew but did not need drug manipulation to prove. It is great that data gathered through drug manipulations converge on those conclusions. Take, for example, one of my recent publications (Doss et al., 2023), which showed that psychedelics can selectively impair memory recollection but spare (or enhance) familiarity even at high doses, while other drugs impaired both processes. This suggests these memory functions are two distinct processes, but there were already lesion populations showing selective impairments of each memory process and loads of fMRI and EEG research that found evidence for their dissociation. But imagine if we only had evidence for this dissociation from drug findings. It would make you wonder if such findings can even be applied to understanding a typical brain. Nonetheless, I do not think pharmacology will ultimately teach me anything new about memory. Of course, I could be wrong, and I am still open to that possibility. Another thing I am interested in is the impact of drugs on memory reconsolidation, a process where retrieved memories become destabilized and altered before being re-stored in a modified state, particularly due to the clinical implications of memory alteration. Memory reconsolidation allows for the updating and integration of new information into existing memories. Research has shown that when people retrieve memories, their mental state becomes very susceptible to external manipulation, such that it is possible to delete specific memories (Haubrich \& Nader, 2018). But it is very difficult to trigger memory reconsolidation processes in humans and much easier in lab rats. One reason for this might be that lab rats have a very homogeneous life, and do not have a lot of other memories. Any type of stimulus used for the induction of new memories, for instance, a loud buzzing sound associated with an electric shock, is likely to be the first time it encounters that sound. I also ALIUS Bulletin n6 (2024) aliusresearch.org/bulletin Doss - On The Cinco Sins of Psychedelic Research wonder whether it would be different in a rat population from the New York subway where they are quite used to such sounds. But in lab rats who are not used to such stimuli, it is very simple to trigger the reconsolidation of fear by re-administering the same sound. By contrast, humans have vastly more associations with any type of stimulus, which makes it very difficult to trigger memory reconsolidation processes in a targeted manner. I would be interested in whether psychedelics (or any other type of drugs) could strengthen or destabilize memories in some manner that makes memory reconsolidation easier. In your critique of research on the subjective effects of psychedelics, you suggest that many researchers are primarily motivated to confirm their own positive experiences with drugs are true. You have criticized the belief in the psychedelic research field that these drugs reveal new insights into brain or mind structure, stating that research has mainly shown how these substances affect the mind without discovering any new underlying principles (Love, 2019. You also noted that scientists privately admit to using psychedelics but avoid public disclosure due to legal issues and fear of losing funding. This seems to paint a picture, that the research is biased because researchers do not disclose their personal experiences with drug use. However, looking at it from another angle, those researchers might argue that their personal experiences can improve the quality of their work. Even though they have personal experiences, their research findings are still based on empirical evidence and do not only reflect their own experiences, even if that was the motivation to research it. This situation might be similar to researchers who decide to study mental health issues like depression because they or their family members have gone through it. It would be interesting to know what drove you to research psychedelics in the first place. When it comes to drugs in general, I had some curiosity since high school. I used to hang around skater kids and participate in rollerblading doing tricks and flips. It was a countercultural scene, and I recall encountering weed around seventh or eighth grade. Initially, I strongly opposed it and believed that skating was all I needed. However, I noticed friends quitting skating for drugs, which made me question why people chose that path. It was interesting yet concerning to see the negative consequences some of my friends experienced. The first time I seriously took notice of hallucinogens was when a friend from high school took dextromethorphan and later tragically took his ALIUS Bulletin n6 (2024) aliusresearch.org/bulletin Doss - On The Cinco Sins of Psychedelic Research own life. This incident made me reflect on the bizarre effects of hallucinogens, such as delusional thinking. Although I was initially against drugs, I found myself strangely fascinated by them, similar to a morbid fear and fascination towards snakes. Over time, my perspective changed, and I realized that people should have the freedom to make their choices as long as they don't harm others. I also recognize that every class of drugs serves a medical purpose, including ketamine, dextromethorphan, and opioids. While some drugs are safer than others, psychedelics seem to be the exception at the moment, but I think it is essential to explore and compare different psychoactive drugs to gain a comprehensive understanding of their unique effects and potential therapeutic uses. Thank you for sharing this history. Following up on my previous question, could you elaborate on why you think first-hand experiences lead to biases and bad science? Could you give some examples? I agree that for some researchers, their initial experience with drugs is what led them down the path of studying them, and this is largely independent of whether or not it taught them anything about how the brain, mind, or even the drugs themselves function. I can relate this to a concussion I had from a skating accident, which did not teach me much about memory, but the experience got me interested in studying it. The main problem is perhaps that we really cannot discuss our personal experiences openly. But on the other hand, I still have the impression that a lot of people enter the field primarily because of their past drug use, and not because they have a deeper scientific interest in neuroscience or psychology. This is similar to how many people want to become neuroscientists because putting pretty brain images makes the science more convincing (McCabe \& Castel, 2008). [M]any studies seem biased, because they are only designed to confirm predictable outcomes. They create the illusion of progress, but I am not convinced we are learning anything new by confirming that psychedelics are psychedelic. My critique of the psychedelic field is largely due to how the nature of the work creates the impression of bias. I hesitate to label it as 'bad science'. It ALIUS Bulletin n6 (2024) aliusresearch.org/bulletin Doss - On The Cinco Sins of Psychedelic Research is not intended as a criticism of individual researchers, but rather the approach the psychedelic field takes as a whole. In this respect, I think some biases affect researchers' perspectives and lead to less insightful study designs. As I have noted before, many studies seem biased, because they are only designed to confirm predictable outcomes. They create the illusion of progress, but I am not convinced we are learning anything new by confirming that psychedelics are psychedelic. There are various examples of this in recent research. An obvious example is a recent study that reported improvements in racial trauma and mental health symptoms following psychedelic experiences (Williams et al., 2021) and they recruited only participants who attributed improvements in racial trauma to psychedelics, already before data collection. This would be very similar to conducting a study where we recruit people to treat their depression with psilocybin, and then conclude that psilocybin-seeking people have depression. It is a circular form of reasoning to report that psychedelics reduce racial trauma, in people who report improvements in racial trauma after taking psychedelics, and I think the primary motivation of this type of research is to prove something that the researchers already believe is true. Another less obvious example is research that shows that LSD makes people react more emotionally in response to music (Kaelen et al., 2015). These results seem just as likely replicated using other psychoactive substances. Just name a recreational drug under the sun that is not used at music concerts. Someone recently told me they take Adderall at concerts, and this caught me off guard, even though it later occurred to me this is probably not that much different from taking cocaine at a music event. I think there is a clear design flaw in claiming that LSD makes music more emotional while not comparing other substances. A balanced approach may be to play different types of music under the influence of different substances. If this type of research showed us the result, for instance, that people started liking electronic music after taking psychedelics, that would be a much more remarkable effect. Finally, I also bring up the study that showed that microdosing LSD dilates the perception of time perception (Yanakieva et al., 2019). The study had participants reproduce the duration of a stimulus (a blue circle) after viewing it. Under LSD, participants tended to reproduce longer intervals. However, this does not prove that psychedelics dilate time perception. If time perception were truly dilated, both the viewing and reproduction phases would be affected equally. This situation is akin to the El Greco fallacy, where it was wrongly assumed that El Greco's elongated painting style was due to a vision problem, not considering that this would have also altered his perception of the canvas. Interestingly, under LSD, they found that participants' estimates of time ALIUS Bulletin n6 (2024) aliusresearch.org/bulletin Doss - On The Cinco Sins of Psychedelic Research duration became more accurate, suggesting an improved alignment of perceived time with memory. The study's design conflated the encoding and retrieval phases of memory, making it unclear whether LSD affects the encoding of time duration or the retrieval process. A study of how these phases are affected separately might be more insightful. Additionally, earlier research (Wackermann et al., 2008) indicated that high and low doses of psilocybin might have the opposite effect on time perception, which raises questions about whether the emphasis on encoding vs. reproduction might shift under different psychedelic doses. Going back to my main point though, I think this case demonstrates how the design and implementation of a study would benefit from expertise in cognitive psychology. Overall, I think such biases should be acknowledged, and we all need to be more critical of each other. I remain skeptical, as I have yet to hear someone explain how their drug experiences have provided insights into the workings of the mind or how drugs work. I think personal experiences may be helpful in examining some basic low-level perceptual phenomena such as color constancy. A researcher skilled in perceptual studies ... might gain valuable insights from a psychedelic experience. Could we entertain the opposite idea, the possibility that having personal experience with psychedelics might make research better? Do you think there are situations where a researcher's own experiences with these drugs could help shape their research in a way that leads to important discoveries or ideas? I think personal experiences may be helpful in examining some basic low-level perceptual phenomena such as color constancy. A researcher skilled in perceptual studies, particularly knowledgeable about the complexities of the magnocellular or parvocellular pathways, might gain valuable insights from a psychedelic experience. Descriptions of psychedelic experiences often suggest they disrupt color constancy, yet the narrators lack the specific vocabulary to pinpoint it. When I initially discovered this concept in psychophysics, I recognized that it might clarify certain psychedelic effects. An expert in visual perception and color constancy could potentially discern the underlying cause of these alterations. Should they experiment with psychedelic substances without any preconceptions, they might question: why does a wall appear purple? Why do violet tones seem more prominent than ALIUS Bulletin n6 (2024) aliusresearch.org/bulletin Doss - On The Cinco Sins of Psychedelic Research greens and reds? Understanding that a white wall reflects a spectrum of colors, they might ponder why psychedelics accentuate the violets. Typically, the perception of specific colors hinges on external factors like the lighting conditions or surrounding colors. Such a researcher might examine these processes and inquire if these are typical characteristics of psychedelic experiences, or if they vary based on the context. For example, if an individual experiences a heightened perception of a wall's violet hues under LSD at a certain time, would the same hues be as pronounced under different circumstances? This field remains largely unexplored, thus a deep comprehension of visual perception combined with a psychedelic experience might offer critical insights for research. Therefore, investigating how psychedelics disrupt certain aspects like constancy, detection of peripheral motion, or other perceptual distortions may be a worthwhile effort. It is conceivable that an individual with extensive knowledge of perception could formulate a theory for the drug's action post-experience. However, the bizarre, extraordinary, or unexplained phenomena reported by users draw more attention. For many people, the value of their psychedelic experiences lies more in the deep insights they gain about their views on life, rather than just changes in basic perception or how they process information. How do you view the scientific efforts to explore these deeper, "high-level" experiences? For example, your colleague David Yaden speaks of the noetic qualities of psychedelic experiences (Yaden et al., 2021; Yaden et al., 2017). Fortier-Davy and Milliere (2020) also note that the heightened 'sense of reality' which is characteristic of psychedelics, is an important factor in distinguishing them from hallucinations induced through anticholinergics, or hallucinations in psychosis-like states. Psychedelic hallucinations seem to be much more lucid compared to anti-cholinergic hallucinations (Fortier, 2018), where people become completely detached from reality and lack basic insights into the fact that they are even hallucinating. What do you think about the study of these phenomena, the current frameworks that investigate psychedelic drug actions in terms of alterations in higher-level, or top-down, aspects of cognition? First of all, people may just as well encounter delusional states when under the influence of psychedelics, such as experiencing complete disconnection from reality with high doses of psilocybin or vivid hallucinations of entities like machine elves with inhaled DMT. Even at lower doses, individuals can develop delusional beliefs or a misleading sense of insight, believing they have gained a profound understanding of the universe or ALIUS Bulletin n6 (2024) aliusresearch.org/bulletin Doss - On The Cinco Sins of Psychedelic Research consciousness that often doesn't translate into their everyday lives. People may report that things feel more real, more important, or more relevant to them, but I maintain that it is crucial to approach these claims with skepticism, akin to how one would view similar reports from individuals with schizophrenia, recognizing that a heightened sense of reality does not equate to actual reality. But for some reason when people report seeing machine elves on DMT, some researchers seem to take the notion that they may exist, seriously. I maintain that it is crucial to approach these claims with skepticism, akin to ... reports from individuals with schizophrenia, recognizing that a heightened sense of reality does not equate to actual reality. But for some reason when people report seeing machine elves on DMT, some researchers seem to take the notion that they may exist, seriously. Regarding your main question, I take a lot of inspiration from Endel Tulving (1985) who introduced the term "noetic consciousness" within the context of memory research. His work also introduced the distinction between episodic and semantic memory, closely linking episodic memory with the concept of autonoetic consciousness and semantic memory with noetic consciousness. Episodic memory allows individuals to mentally travel in time, recalling past experiences or imagining future ones, a process heavily reliant on the hippocampus, and involving a form of consciousness that entails an intact autobiographical self. On the other hand, noetic consciousness, associated with semantic memory, involves a basic sense of knowing or familiarity that operates independently of the hippocampus, enabling recognition without personal experiential recall. For instance, if we show people a list of words, and then disrupt the hippocampus they would still have a vague sense of which items were more or less familiar. If you take the hippocampus offline, then this sense of familiarity can be attached to irrelevant information. This is also in line with research demonstrating how the feeling of insight can be increased via semantic priming manipulation and misattributed to irrelevant ideas (Grimmer et al., 2022; McGovern et al., 2023). Thus it seems worth exploring how psychedelics operate, either by taking the hippocampus offline or through some other mechanism that boosts the sense of familiarity to test whether they create false insights or memories. The notion that psychedelic experiences ALIUS Bulletin n6 (2024) aliusresearch.org/bulletin Doss - On The Cinco Sins of Psychedelic Research provide a heightened sense of reality or true insight is, in my view, questionable, and we should just as well investigate whether they cause misattribution that can then lead to a belief in the reality or importance of these ideas or insights, despite a lack of objective evidence to support them, which is rather delusional. In my research, I am most interested in uncovering new and fascinating aspects of these drugs, for instance, if they can enhance or impair memory, whether they induce sleep or wakefulness, and a variety of other effects. When investigating a drug that produces very strong subjective effects, I think it crucial to compare them with other psychoactive drugs that have similar characteristics. For example, psychedelics increase wakefulness, so why are we not comparing their effects with stimulants? Psychedelics also cause hallucinations, so we could also compare them with other NMDA antagonists like ketamine. But grouping ketamine, and serotonergic hallucinogens in one group as "psychedelics" based on similarities in their subjective effects does not advance our understanding of their effects. What is your take on more ambitious frameworks like the relaxed beliefs under psychedelics (REBUS) model (Carhart-Harris \& Friston, 2019) that try to explain higher-order phenomena within the framework of predictive processing (Wiese \& Metzinger, 2017) or the Free Energy Principle (Friston, 2010)? Do these frameworks provide a solid scientific foundation for exploring how subjective experiences, like the sensation of insight, correspond to mechanisms, such as the reduction of prediction error? I think the REBUS model shares the same shortcomings as predictive processing, which remains unfalsifiable (Friston et al., 2018). It seems like loads of subjective phenomena are thrown onto a prediction error: Oh, it's! It's rewarding! Oh, it produces insight! Oh, it produces surprise. Oh, it produces fear. Just make up your goddamn mind! Speaking to the notion that feelings of insight could be a result of specific mechanistic processes such as elevated rates of prediction error, or priors being updated at an uncommonly high level, the answer is not clear-cut. Various circuits form some kind of prediction error but I don't know how you can map something like that onto various subjective effects, certainly not a single one. I am sure it varies quite a bit by circuit. I am generally skeptical regarding its applicability due to this over-attribution of various subjective phenomena to prediction error. ALIUS Bulletin n6 (2024) aliusresearch.org/bulletin Doss - On The Cinco Sins of Psychedelic Research My other critical issue with the REBUS model is its apparent lack of functional alignment and precision necessary to accurately map and predict the effects of psychedelics on cognition and behavior. The model puts forth the idea that under the influence of psychedelics, high-level cognition should be more impaired than low-level cognition. Yet, psychedelics can disrupt low-level learning forms, such as prepulse inhibition, while potentially enhancing higher-level functions like mnemonic familiarity (Doss et al., 2022; Doss et al., 2023). The REBUS model's claim that the hippocampus is a low-level structure, and the connectivity between the hippocampus and the default mode network increases under psychedelics, does not fully align with observed data. Studies show that impairment of the perirhinal cortex but not the hippocampus leads to impairments of familiarity (Bowles et al., 2007). Consequently, if psychedelics enhance familiarity, it suggests the engagement of high-level processes, not low-level ones as suggested by the model. The way REBUS is written seems to be agnostic to cognitive psychology in many ways, except when it suits its purposes. REBUS, much like the Free Energy Principle, tries to be a theory of everything. REBUS is supposed to be an explanation at the computational level. Until you can show me some very specific behaviors, I don't think that I'm going to buy into it. While the REBUS model's concept of heightened bottom-up prediction error aligns with reports from individuals under the influence of psychedelics - who often describe everything as being novel or surprising - it falls short in explaining other psychedelic experiences, such as deja vu. Some of the things you'd expect from REBUS like greater brain connectivity and greater brain states are found in patients with obsessive-compulsive disorder (OCD). Are they always tripping out then? Even if you take the brain measures themselves, such as entropy, they don't have any specificity. Entropy increases with coffee. There's not a one-to-one mapping between the subjective experience, but it's meant to be somewhat expansive. The framework starts to break down when you begin attributing real functions to these brain predictions. The way REBUS is written seems to be agnostic to cognitive psychology in many ways, except when it suits its purposes. Currently, there's a tendency to generalize findings based on a small number of ALIUS Bulletin n6 (2024) aliusresearch.org/bulletin Doss - On The Cinco Sins of Psychedelic Research participants. Everybody's piggybacking on REBUS instead of challenging it. I think that's what we should be doing - challenging different theories. Before the REBUS model was introduced, Vollenweider and Geyer (2001) had already proposed the idea that psychedelics could disrupt thalamic gating, causing a shift from top-down to bottom-up cognition. Nicholas Langlitz (2007) traces this idea's lineage back to Aldous Huxley's perennial concept of the mind as a "reducing valve" that filters sensory information to create a coherent reality. Huxley suggested that psychedelics expand this valve, allowing for a broader spectrum of perception and consciousness by permitting a more unfiltered influx of sensory data. Considering this backdrop and the hypothesis that psychedelics flatten or disrupt high-level priors, potentially enhancing susceptibility to priming and altering the content of hallucinations, what are your broader thoughts on this framework? How do you view the interplay between top-down and bottom-up processes in this context, and do you believe it offers a valid avenue for experimental exploration? How might this altered susceptibility to priming and the selective disruption of higher-level information affect the nature of hallucinations and individuals' susceptibility to them? What I do not like about the notion of the reducing valve is that it could be interpreted as seeing things that are there, without distortions in your visual field, and the notion that everything is due to a relaxation of priors, like visual distortions. In theory, certain thoughts may be more likely to occur that would normally be filtered out. It would be interesting to explore this hypothesis by comparing hallucinations across cultures, for instance, whether people from Brazil are more likely to hallucinate jaguars on Ayahuasca compared with someone from a Western background. But on the other hand, such hallucinations require high-level representations which are allegedly disrupted by psychedelics. Despite predictions of disrupted high-level representations, such hallucinations continue to appear even on really high doses. People report visions of complex landscapes or encounters with deceased relatives, indicating that high-level concepts remain intact under psychedelics, which challenges the "flattened landscape" metaphor. This raises the question of where the hierarchy begins and ends, of what is considered low-level versus high-level. For instance, on a behavioral level, psilocybin impairs the perception of motion (Carter et al., 2004) whereas alcohol does the opposite (Wang et al., 2018), but in my view, these are still pretty low levels of information processing. Then there are instances of certain high-level functions being facilitated, like semantic priming (Spitzer et al., 1996), which is further upstream than these visual effects. Therefore, I propose that there are likely cases of facilitated top-down ALIUS Bulletin n6 (2024) aliusresearch.org/bulletin Doss - On The Cinco Sins of Psychedelic Research processing under psychedelics, which also do not comply with the reducing valve or thalamic gating model. Looking at it on the circuit level, there is a hierarchy of back-and-forth projections throughout the cortex, and as you move from lower to higher levels of the cortex, there are also corresponding levels in the thalamus receiving and projecting back. The expected impairment of higher-order cognition being more pronounced than lower-order cognition is contradicted by the fact that even basic forms of learning, such as prepulse inhibition, are disrupted by psychedelics. From a subjective standpoint, I would also argue that mnemonic processes such as familiarity are enhanced by psychedelic substances and that these depend on structures such as the perirhinal cortex which is situated at the top of the ventral visual stream. This does not align well with the idea that psychedelics break down or flatten a prediction hierarchy. One of the explicit predictions from the REBUS model is that there should be greater effective connectivity, in other words, greater information flow from the hippocampus to the cortex. But this doesn't make sense considering hippocampus-dependent memory or recollection is impaired by psychedelics. Familiarity or the feeling of knowing is more dependent on the cortex. We tried to make sense of the REBUS model in terms of the circuit implementation. When it comes to the circuits, I think the cortico-striatal thalamo-cortical (CSTC) model (Doss et al., 2021) has the most evidence behind it and it does have some explanatory value for certain phenomena, like impairments in attention and inability to focus on something. In discussions with Phil Corlett, I have contemplated whether the neural dynamics under psychedelics resemble more of a "bumpy seesaw" than a "flattened landscape." In discussions with Phil Corlett, I have contemplated whether the neural dynamics under psychedelics resemble more of a "bumpy seesaw" than a "flattened landscape." This "bumpy seesaw" metaphor, considering the slope as the flow of information within the hierarchy, might more accurately reflect the interplay between top-down and bottom-up information processing. While a flattened landscape suggests a uniform distribution of influence, the seesaw model implies a dynamic balance, with information oscillating between higher and lower processing levels. This back-and-forth motion suggests that both bottom-up and top-down information can penetrate more readily under the influence of psychedelics, challenging the notion of a completely flattened ALIUS Bulletin n6 (2024) aliusresearch.org/bulletin Doss - On The Cinco Sins of Psychedelic Research hierarchy. Essentially, this interaction includes a complex network of projections within the cortex and between the cortex and the thalamus, allowing for a nuanced exchange of information across different levels of cognition. Can you outline your main critique of studies aiming to establish a link between the subjective effects of drugs and resting-state brain connectivity? Numerous studies have attempted to establish a link between the subjective effects of drugs and brain activity during periods of rest, but a key issue is the lack of replication in the observed correlations. For instance, while one study by Carhart-Harris et al. (2016) found that decreases in default mode network connectivity correlate with ego dissolution, another study by Muller et al. (2018) failed to replicate this finding despite controlling for head motion and having more subjects. This lack of consistency highlights the challenge of establishing a one-to-one mapping between brain activity and subjective drug effects. My takeaway from this, is researchers must constrain their search space and exert more experimental control to gain a better understanding of the effects of drugs on the mind and brain. Although differences in retrospective reports, behavior, and brain activity may offer some insights, a one-to-one mapping between brain activity and subjective experience is not yet achievable. Even with advanced neuroimaging techniques such as fMRI and EEG, diagnosing depression remains a challenge. One potential approach to narrowing down the search space is to focus on behavioral measures, but the interpretability of the results may be limited, and you cannot just widely speculate on pet theories about the brain based on such behavioral measures. One of the things I have suggested is that all of the resting state analysis should be happening on data gathered while participants are completing a task. Some proponents of resting-state activity may argue that the drug's effects on intrinsic connectivity should manifest across any task manipulation. But in my view, if an effect is too weak to persist when stimuli are introduced, it is unlikely to be a result of drug-induced changes in intrinsic connectivity. The consistent observation of decreased resting-state connectivity of the default mode network across various drugs suggests that this effect may not be as interesting as previously thought. This effect may be related to the fact that the default mode network is more active during internal mental processes, such as representing the world internally when attention is not directed outward. However, under the influence of drugs, attention may be periodically redirected to the external environment, even with alcohol. This may cause ALIUS Bulletin n6 (2024) aliusresearch.org/bulletin Doss - On The Cinco Sins of Psychedelic Research brain activity similar to those observed with stimulants and psychedelics. Therefore, failing to control for these external factors may lead to inconsistent results, as seen with various ketamine trials: three studies showed decreases in default mode network connectivity (Bonhomme et al., 2016; Scheidegger et al., 2012; Zacharias et al., 2020), two studies showed no effect (Mueller et al., 2018; Niesters et al., 2012), and the study with the largest sample size of fifty-three participants found increased default mode network connectivity. For me, these mixed results indicate a null effect. Has the enthusiasm for the Default Mode Network (DMN) in psychedelic research mirrored the broader hype in Cognitive Neuroscience and Psychiatry for its potential diagnostic value? Many researchers viewed the Default Mode Network as one of the most reliable markers of depression (e.g. Scalabrini et al., 2020). Would you balance your critique without entirely dismissing the DMN's relevance as a therapeutic marker or the study of cognition? By focusing solely on resting-state activity, researchers can manipulate and interpret brain data using unlimited mathematical techniques to support any conclusion they wish to make. From my perspective, the exclusive investigation of the brain in a resting state is not truly representative of Cognitive Neuroscience because it fails to consider cognition. The association between depression and Default Mode Network (DMN) connectivity is also questionable due to inconsistent findings. While the meta-analysis by Kaiser et al. (2015) found that patients with depression exhibit increased DMN connectivity, a study by Yan et al. (2019) observed reduced connectivity across a 1300-participant sample. Although methodological differences could account for the discrepancy, the main issue seems more fundamental to me. By focusing solely on resting-state activity, researchers can manipulate and interpret brain data using unlimited mathematical techniques to support any conclusion they wish to make. Russel Poldrack (2011) previously pointed out the problem of reverse inference in non-resting state data, where brain activity during a resting state is used to make assumptions about the mind rather than the other way around. Moreover, I believe this approach discredits the field of experimental psychology by ignoring behavior. But there is plenty of blame to share, as it is reasonable to suggest that some cognitive psychologists aim to enter the field ALIUS Bulletin n6 (2024) aliusresearch.org/bulletin Doss - On The Cinco Sins of Psychedelic Research of neuroscience in this manner, all while scrutinizing clinical psychology for lacking enough rigor. How do you respond to the argument that resting-state assessments provide a more naturalistic evaluation of brain activity under the influence of drugs, especially considering these substances can lead to significant disengagement from cognitive tasks and cause substantial fluctuations in attention, making standard behavioral methods challenging to apply? There are a few different aspects to consider here, but I think there is a general reluctance among many researchers to use behavioral tasks, due to a prevailing attitude that participants, under the influence, lack engagement and interest in the task. I generally disagree with this view. However, it is essential to acknowledge the limitations of relying exclusively on standardized methods. Cognitive batteries in psychedelics studies, such as those by Pokorny et al. (2020), often indicate cognitive impairments, but these results may not be very informative on their own. Instead, a more in-depth examination of specific tasks related to attention, memory, and other cognitive processes is necessary, as well as studies that directly compare different drugs, like dextromethorphan and psilocybin. For example, Barrett et al. (2018), reveal intriguing distinctions: dextromethorphan significantly impairs episodic memory encoding with little effect on working memory, whereas even the lowest dose of psilocybin impairs working memory but minimally affects episodic memory. These findings challenge the idea of generalized disengagement and highlight the need for a nuanced understanding of the cognitive effects of psychedelics. If task engagement drops significantly, it becomes evident through participants' responses, and monitoring reaction times and response accuracy can ensure continued engagement. We must delve deeper. For instance, attention is not a single concept; psychedelics may enhance certain aspects of attention. Steve Luck demonstrated that individuals with lower working memory capacity sometimes remember non-cued stimuli better than those with higher capacities (Luck et al., 2002), due to fluctuating attention that increases the likelihood of including the non-cued stimulus in their mental workspace. Although this may impair overall working memory, it can lead to better performance in some cases. Investigating this effect with psychedelics, especially serotonergic ones as opposed to substances like dextromethorphan, could be intriguing. Attention encompasses various types and levels, such as covert versus overt, or bottom-up versus top-down. A study by Spitzer et al. (2001) showed that the pop-out effect, a low-level visual effect considered a type of attention, was enhanced by the R-enantiomer of ALIUS Bulletin n6 (2024) aliusresearch.org/bulletin Doss - On The Cinco Sins of Psychedelic Research methylenedioxyethylamphetamine (MDEA) but not the S-enantiomer, suggesting a specific type of attention enhancement. Although this finding requires replication, it opens new research avenues. When examining episodic memory impairment, the issue is complex. A deeper look into the data suggests that cortical-dependent memory encoding might be enhanced under certain conditions. In summary, to understand the peculiar phenomena where both high novelty and familiarity occur under the effects of psychedelics, we need to look beyond the surface and move away from overarching theories that explain little. How do you view the use of experience sampling methods during resting-state (Fell, 2013; Gonzalez-Castillo et al., 2021; Heimann et al.) to explore the mapping of these experiences more intricately via a more fine-grained qualitative methodology? Utilizing experience sampling techniques may probably produce better data, but you would probably need to probe people about 100 times an hour to get reliable samples to see what is going on. It also introduces the problem that participants are going to be expecting the next probe, so I think showing participants a standardized stimulus, such as a movie, would already be better. I still think there is something to be said about getting just good behavioral results to make an educated reverse inference of how people can still perform certain cognitive functions under the influence of a drug, but perhaps by using the brain in a different way than usual. Your critique of psychedelic research appears to focus on skepticism towards subjective effect measurements and a preference for more objective behavioral paradigms, reflecting your background in cognitive psychology. You argue that relying on subjective questionnaires tends to validate preconceived notions, a practice you have (in public lectures) termed as conducting "Duh"-science. But considering the tradition in psychology to start with folk psychological understandings before advancing to scientifically validated measures, is psychedelic research still navigating early stages, where testing intuitive notions about drug effects is a necessary step? My problem with current research is that people are not even testing folk intuitions, but rather coming in with biases from their drug-using experiences. I am sick of people acting as if they have never experienced the drug's effects. However, when you look at the things that people choose to measure, it is clear whether or not a researcher has tried certain drugs. ALIUS Bulletin n6 (2024) aliusresearch.org/bulletin Doss - On The Cinco Sins of Psychedelic Research Investigating obvious subjective effects, such as mystical experiences or whether music sounds better under the influence of drugs, makes it clear to participants how the researchers expect them to respond. Various cultural tropes will also lead them to respond in this manner. That is why I have more faith in cognitive psychology because researchers can set up an experiment in which they stand to lose sometimes, but if they can still show an effect, it makes it all the more interesting. Based on the Mystical Experience Questionnaire scores, I would bet that many recreational drugs could induce mystical experiences at high doses. Perhaps it is useful to confirm these facts, but would you expect anything different? Despite the 'psychedelic renaissance', the reliance on predictable subjective measures leaves me questioning what new insights, if any, we have gained about the psychological impacts of these substances. I have definitely not gained any new insights about the mind by researching [psychedelic] drugs. ... what have you learned? I am guessing not much. I have definitely not gained any new insights about the mind by researching these drugs. So let me ask you, what have you learned? I am guessing not much. I will aim to provide a diplomatic answer to that question. Although I largely agree with all your methodological criticisms. particularly because many of your concerns seem reminiscent of my favorite term, "Neurophrenology" (or blobology), which involves trying to directly map mental phenomena to specific brain regions (Anderson, 2021; Friston, 2002; Huskey, 2016; Kosik, 2003). At the same time, I would still like to counter the polemical nature of the statement. Firstly, because scientific discoveries rarely occur in "Eureka" moments, but rather through incremental puzzle-solving, and I would not claim that psychedelic research is exempt from this or holds any kind of exceptional, paradigm-shifting status. However, from a historical perspective, I do think that psychedelics have significantly contributed to the modern understanding of the brain, at least regarding its neurochemical basis. David Nichols (2013) articulates this better than I can, explaining that research on LSD played a significant role in elucidating the role of serotonin in the brain, leading to a deeper understanding of the neurochemical basis of the mind. The discovery of serotonin's function in the mind laid the foundation for a whole new field of psychopharmacology aimed at ALIUS Bulletin n6 (2024) aliusresearch.org/bulletin Doss - On The Cinco Sins of Psychedelic Research elucidating the neurochemical basis of the mind. Moreover, this research helped alleviate social stigma towards parents of children with autism or schizophrenia, countering the blame placed on refrigerator mothers' who were accused of withholding emotions from their children. So, in that sense, psychedelic research has already revealed a lot about the mind, albeit perhaps not through the current renaissance. Moreover, I still think that Huxley's description of psychedelics as 'mind-revealing' has shaped current research in a meaningful way and that both the hypotheses around thalamo-cortical gating (Vollenweider \& Geyer, 2001) and the relaxation of top-down constraints (Carhart-Harris \& Friston, 2019) present a genuine attempt to reinterpret the notion of psychedelics as mind-revealing in a way that is conducive to scientific investigation and testing. Even though I acknowledge the caveats of the REBUS model, I argue that the scientific disagreements it provokes are valuable because they allow us to conduct experiments and build theoretical models that become more precise over time. Similar to how psychology begins by verifying folk intuitions, I think we can start with a broader theory and refine the details as we advance our understanding through encountering discrepancies. Furthermore, I would like to advocate for the exploration of bizarre and profound subjective phenomena, as illustrated by the historical example of Hans Berger. His investigation into electroencephalography (EEG) was motivated by a near-death experience that led him to believe he had experienced telepathy (La Vaque, 1999; Millett, 2001). Despite his initial aim to explain the neurophysiological basis of telepathy, Berger instead discovered alpha waves. I believe this principle applies to curiosity about strange phenomena, such as hallucinations of entities, in that they may reveal intriguing research avenues, provided the research is conducted rigorously. Despite my agreement with many of your criticisms regarding the use of subjective measures in current research, I believe entirely dismissing subjective data would be a mistake. Valuable insights can be gleaned from comparing the phenomenology of different states of consciousness, thereby enriching our understanding of cognition. For example, examining the qualitative aspects of various hallucinations can serve as a valid starting point for identifying differences in their cognitive and neurochemical underpinnings. Furthermore, your critique of subjective effects resonates with a long-standing debate within cognitive psychology about the reliability of introspective reports. Jack and Roepstorff (2002, 2003) emphasize that introspective evidence, crucial for understanding consciousness and subjective states, informs all aspects of cognitive science, including research and hypothesis generation. They advocate for its integration with behavioral and physiological data to gain comprehensive insights into cognitive processes. In the same vein, I challenge the notion of conducting ALIUS Bulletin n6 (2024) aliusresearch.org/bulletin Doss - On The Cinco Sins of Psychedelic Research cognitive science or psychedelic research entirely without introspection, partly inspired by discussions on 'cognitive phenomenology'-the idea that certain cognitive processes possess a phenomenal quality. Addressing phenomena such as the intense feelings of insight reported by people under the influence of psychedelics can serve as a basis for investigating corresponding cognitive mechanisms. Therefore, I advocate for an approach where cognitive and subjective measures mutually inform one another. This brings me to my final question: Do you still remain critical of the idea that understanding experience from a first-person perspective, when approached with reflective intention, could enrich scientific practice? In my view, your first point about serotonin is more about the brain than the mind. While I think there's something to that, take, for example, the discovery of 5-HT2A receptors on glial cells. That by itself wouldn't tell us much (yet) about the mind. Regarding Huxley, I am not sure what the term "mind-revealing" has to do with thalamic gating or REBUS (though his "reducing valve" analogy certainly alludes to the former). I think research based initial intuition is fine, but we can not forget the plethora of foundational knowledge in psychology and neuroscience, which is why I tend to like the thalamic gating model, as it respects the known circuitry of the cortex, basal ganglia, and thalamus (though it could use more behavioral tests and refinement in terms of incorporating efferent copies). To be clear, I have never claimed to be against introspection, and I am pretty unimpressed by traditional behaviorism. Motivation by personal interest or experience can be a powerful tool for answering research questions, but we always have to remain aware that such motivation can also drive us to negate answers we do not like. To be clear, I have never claimed to be against introspection, and I am pretty unimpressed by traditional behaviorism. In fact, most cognitive measures typically rely on some form of introspection. What I find uninformative are studies that simply aim to find some obvious subjective measure (e.g., THC makes people hungry), as they are clearly prone to confirmation bias. I think psychedelic research in general needs stronger inference, and to make subjective measures more informative. The framing of questionnaires needs to be changed up in an effort to either refute or strengthen an effect across multiple experiments in a single ALIUS Bulletin n6 (2024) aliusresearch.org/bulletin Doss - On The Cinco Sins of Psychedelic Research publication, similar to how studies in cognitive psychology might contain multiple experiments to confirm an effect. But salami slice or perish I suppose. References: Anderson, M. L. (2021). After phrenology: Neural reuse and the interactive brain. MIT Press. Barrett, F. S., Carbonaro, T. M., Hurwitz, E., Johnson, M. W., \& Griffiths, R. R. (2018). Double-blind comparison of the two hallucinogens psilocybin and dextromethorphan: effects on cognition. Psychopharmacology, 235, 2915-2927. Bonhomme, V., Vanhaudenhuyse, A., Demertzi, A., Bruno, M.-A., Jaquet, O., Bahri, M. A., Plenevaux, A., Boly, M., Boveroux, P., Soddu, A., Brichant, J. F., Maquet, P., \& Laureys, S. (2016). Resting-state Network-specific Breakdown of Functional Connectivity during Ketamine Alteration of Consciousness in Volunteers. Anesthesiology, 125(5), 873-888. https://doi.org/10.1097/aln.0000000000001275 Bowles, B., Crupi, C., Mirsattari, S. M., Pigott, S. E., Parrent, A. G., Pruessner, J. C., Yonelinas, A. P., \& Kohler, S. (2007). Impaired familiarity with preserved recollection after anterior temporal-lobe resection that spares the hippocampus. Proceedings of the National Academy of Sciences, 104(41), 16382-16387. Carhart-Harris, R. L., \& Friston, K. J. (2019). REBUS and the anarchic brain: toward a unified model of the brain action of psychedelics. Pharmacological reviews, 71(3), 316-344. Carhart-Harris, R. L., Muthukumaraswamy, S., Roseman, L., Kaelen, M., Droog, W., Murphy, K., Tagliazucchi, E., Schenberg, E. E., Nest, T., \& Orban, C. (2016). Neural correlates of the LSD experience revealed by multimodal neuroimaging. Proceedings of the National Academy of Sciences, 113(17), 4853-4858. Carter, O. L., Pettigrew, J. D., Burr, D. C., Alais, D., Hasler, F., \& Vollenweider, F. X. (2004). Psilocybin impairs high-level but not low-level motion perception. Neuroreport, 15(12), 1947-1951. https://doi.org/10.1097/00001756-200408260-00023 Doss, M. K., Samaha, J., Barrett, F. S., Griffiths, R. R., de Wit, H., Gallo, D. A., \& Koen, J. D. (2023). Unique effects of sedatives, dissociatives, psychedelics, stimulants, and cannabinoids on episodic memory: A review and reanalysis of ALIUS Bulletin n6 (2024) aliusresearch.org/bulletin Doss - On The Cinco Sins of Psychedelic Research acute drug effects on recollection, familiarity, and metamemory. Psychological Review. Doss, M. K., Madden, M. B., Gaddis, A., Nebel, M. B., Griffiths, R. R., Mathur, B. N., \& Barrett, F. S. (2021). Models of psychedelic drug action: modulation of cortical-subcortical circuits. Brain, 145(2), 441-456. https://doi.org/10.1093/brain/awab406 Doss, M. K., Samaha, J., Barrett, F. S., Griffiths, R. R., de Wit, H., Gallo, D. A., \& Koen, J. D. (2023). Unique effects of sedatives, dissociatives, psychedelics, stimulants, and cannabinoids on episodic memory: A review and reanalysis of acute drug effects on recollection, familiarity, and metamemory. Psychological Review. Doss, M. K., Weafer, J., Gallo, D. A., \& de Wit, H. (2018). (9)-Tetrahydrocannabinol at Retrieval Drives False Recollection of Neutral and Emotional Memories. Biol Psychiatry, 84(10), 743-750. https://doi.org/10.1016/j.biopsych.2018.04.020 Fell, J. (2013). Unraveling inner experiences during resting state. In (Vol. 7, pp. 409): Frontiers Media SA. Fortier-Davy, M., \& Milliere, R. (2020). The multi-dimensional approach to drug-induced states: a commentary on Bayne and Carter's "dimensions of consciousness and the psychedelic state". Neuroscience of Consciousness, 2020(1), niaa004. Fortier, M. (2018). Sense of reality, metacognition, and culture in schizophrenic and drug. Metacognitive diversity: An interdisciplinary approach, 343. Friston, K. (2002). Beyond phrenology: what can neuroimaging tell us about distributed circuitry? Annual review of neuroscience, 25(1), 221-250. Friston, K. (2010). The free-energy principle: a unified brain theory? Nature Reviews Neuroscience, 11(2), 127-138. Friston, K., Fortier, M., \& Friedman, D. (2018). Of woodlice and men: A Bayesian account of cognition, life and consciousness. An interview with Karl Friston. ALIUS Bulletin, 2, 17-43. Gonzalez-Castillo, J., Kam, J. W., Hoy, C. W., \& Bandettini, P. A. (2021). How to interpret resting-state fMRI: ask your participants. Journal of Neuroscience, 41(6), 1130-1141. ALIUS Bulletin n6 (2024) aliusresearch.org/bulletin Doss - On The Cinco Sins of Psychedelic Research Grimmer, H., Laukkonen, R., Tangen, J., \& von Hippel, W. (2022). Eliciting false insights with semantic priming. Psychonomic Bulletin \& Review, 29(3), 954-970. https://doi.org/10.3758/s13423-021-02049-x Haubrich, J., \& Nader, K. (2018). Memory Reconsolidation. In R. E. Clark \& S. J. Martin (Eds.), Behavioral Neuroscience of Learning and Memory (pp. 151-176). Springer International Publishing. https://doi.org/10.1007/7854\_2016\_463 Heimann, H., Pletscher, A., \& Ladewig, D. Experience of time and space in model psychoses. In. Huskey, R. (2016). Beyond blobology: Using psychophysiological interaction analyses to investigate the neural basis of human communication phenomena. Innovative methods in media and communication research, 123-138. Kaelen, M., Barrett, F. S., Roseman, L., Lorenz, R., Family, N., Bolstridge, M., Curran, H. V., Feilding, A., Nutt, D. J., \& Carhart-Harris, R. L. (2015). LSD enhances the emotional response to music. Psychopharmacology, 232, 3607-3614. Kaiser, R. H., Andrews-Hanna, J. R., Wager, T. D., \& Pizzagalli, D. A. (2015). Large-Scale Network Dysfunction in Major Depressive Disorder: A Meta-analysis of Resting-State Functional Connectivity. JAMA Psychiatry, 72(6), 603-611. https://doi.org/10.1001/jamapsychiatry.2015.0071 Kosik, K. S. (2003). Beyond phrenology, at last. Nature Reviews Neuroscience, 4(3), 234-239. La Vaque, T. (1999). The history of EEG hans berger: psychophysiologist. A historical vignette. Journal of Neurotherapy, 3(2), 1-9. Love, S. (2019). Should Psychedelic Researchers Trip on the Drugs They're Studying? VICE. https://www.vice.com/en/article/qv7jkx/the-ethics-of-taking-the-drugs-you-stu dy Luck, S. J., Woodman, G. F., Schmidt, B. K., Vogel, E. K., \& Vecera, S. P. (2002). The effects of attentional capture on visual working memory. Journal of Vision, 2(7), 272-272. McCabe, D. P., \& Castel, A. D. (2008). Seeing is believing: The effect of brain images on judgments of scientific reasoning. Cognition, 107(1), 343-352. ALIUS Bulletin n6 (2024) aliusresearch.org/bulletin Doss - On The Cinco Sins of Psychedelic Research McGovern, H., Grimmer, H., Doss, M., Hutchinson, B., Timmermann, C., Lyon, A., Corlett, P. R., \& Laukkonen, R. (2023). The power of insight: How psychedelics may solicit false beliefs. Millett, D. (2001). Hans Berger: From psychic energy to the EEG. Perspectives in biology and medicine, 44(4), 522-542. Morris, H. (2022). POD 52: Cinco Sins of Psychedelic Science with Dr. Manoj Doss. In. Mueller, F., Musso, F., London, M., de Boer, P., Zacharias, N., \& Winterer, G. (2018). Pharmacological fMRI: Effects of subanesthetic ketamine on resting-state functional connectivity in the default mode network, salience network, dorsal attention network and executive control network. NeuroImage: Clinical, 19, 745-757. https://doi.org/https://doi.org/10.1016/j.nicl.2018.05.037 Muller, F., Dolder, P. C., Schmidt, A., Liechti, M. E., \& Borgwardt, S. (2018). Altered network hub connectivity after acute LSD administration. NeuroImage: Clinical, 18, 694-701. https://doi.org/https://doi.org/10.1016/j.nicl.2018.03.005 Nichols, D. E. (2013). Serotonin, and the past and future of LSD. MAPS Bull, 23(1), 20-23. Niesters, M., Khalili-Mahani, N., Martini, C., Aarts, L., van Gerven, J., van Buchem, Mark A., Dahan, A., \& Rombouts, S. (2012). Effect of Subanesthetic Ketamine on Intrinsic Functional Brain Connectivity: A Placebo-controlled Functional Magnetic Resonance Imaging Study in Healthy Male Volunteers. Anesthesiology, 117(4), 868-877. https://doi.org/10.1097/ALN.0b013e31826a0db3 Pokorny, T., Duerler, P., Seifritz, E., Vollenweider, F. X., \& Preller, K. H. (2020). LSD acutely impairs working memory, executive functions, and cognitive flexibility, but not risk-based decision-making. Psychological medicine, 50(13), 2255-2264. Poldrack, R. A. (2011). Inferring mental states from neuroimaging data: from reverse inference to large-scale decoding. Neuron, 72(5), 692-697. Scalabrini, A., Vai, B., Poletti, S., Damiani, S., Mucci, C., Colombo, C., Zanardi, R., Benedetti, F., \& Northoff, G. (2020). All roads lead to the default-mode network-global source of DMN abnormalities in major depressive disorder. Neuropsychopharmacology, 45(12), 2058-2069. https://doi.org/10.1038/s41386-020-0785-x ALIUS Bulletin n6 (2024) aliusresearch.org/bulletin Doss - On The Cinco Sins of Psychedelic Research Scheidegger, M., Walter, M., Lehmann, M., Metzger, C., Grimm, S., Boeker, H., Boesiger, P., Henning, A., \& Seifritz, E. (2012). Ketamine decreases resting state functional network connectivity in healthy subjects: implications for antidepressant drug action. Spitzer, M., Franke, B., Walter, H., Buechler, J., Wunderlich, A. P., Schwab, M., Kovar, K.-A., Hermle, L., \& Gron, G. (2001). Enantio-selective cognitive and brain activation effects of N-ethyl-3, 4-methylenedioxyamphetamine in humans. Neuropharmacology, 41(2), 263-271. Spitzer, M., Thimm, M., Hermle, L., Holzmann, P., Kovar, K.-A., Heimann, H., Gouzoulis-Mayfrank, E., Kischka, U., \& Schneider, F. (1996). Increased activation of indirect semantic associations under psilocybin. Biological psychiatry, 39(12), 1055-1057. Tulving, E. (1985). Memory and consciousness. Canadian Psychology/Psychologie canadienne, 26(1), 1. Vollenweider, F. X., \& Geyer, M. A. (2001). A systems model of altered consciousness: integrating natural and drug-induced psychoses. Brain research bulletin, 56(5), 495-507. Wackermann, J., Wittmann, M., Hasler, F., \& Vollenweider, F. X. (2008). Effects of varied doses of psilocybin on time interval reproduction in human subjects. Neuroscience Letters, 435(1), 51-55. https://doi.org/https://doi.org/10.1016/j.neulet.2008.02.006 Wang, Z., Wang, H., Tzvetanov, T., \& Zhou, Y. (2018). Moderate acute alcohol intoxication increases visual motion repulsion. Scientific Reports, 8(1), 1607. https://doi.org/10.1038/s41598-018-19932-8 Wiese, W., \& Metzinger, T. (2017). Vanilla PP for philosophers: A primer on predictive processing. Williams, M. T., Davis, A. K., Xin, Y., Sepeda, N. D., Grigas, P. C., Sinnott, S., \& Haeny, A. M. (2021). People of color in North America report improvements in racial trauma and mental health symptoms following psychedelic experiences. Drugs: Education, Prevention and Policy, 28(3), 215-226. https://doi.org/10.1080/09687637.2020.1854688 Yaden, D. B., Johnson, M. W., Griffiths, R. R., Doss, M. K., Garcia-Romeu, A., Nayak, S., Gukasyan, N., Mathur, B. N., \& Barrett, F. S. (2021). Psychedelics and ALIUS Bulletin n6 (2024) aliusresearch.org/bulletin Doss - On The Cinco Sins of Psychedelic Research consciousness: Distinctions, demarcations, and opportunities. International Journal of Neuropsychopharmacology, 24(8), 615-623. Yaden, D. B., Le Nguyen, K. D., Kern, M. L., Wintering, N. A., Eichstaedt, J. C., Schwartz, H. A., Buffone, A. E., Smith, L. K., Waldman, M. R., \& Hood Jr, R. W. (2017). The noetic quality: A multimethod exploratory study. Psychology of Consciousness: Theory, Research, and Practice, 4(1), 54. Yan, C.-G., Chen, X., Li, L., Castellanos, F. X., Bai, T.-J., Bo, Q.-J., Cao, J., Chen, G.-M., Chen, N.-X., \& Chen, W. (2019). Reduced default mode network functional connectivity in patients with recurrent major depressive disorder. Proceedings of the National Academy of Sciences, 116(18), 9078-9083. Yanakieva, S., Polychroni, N., Family, N., Williams, L. T. J., Luke, D. P., \& Terhune, D. B. (2019). The effects of microdose LSD on time perception: a randomised, double-blind, placebo-controlled trial. Psychopharmacology, 236(4), 1159-1170. https://doi.org/10.1007/s00213-018-5119-x Zacharias, N., Musso, F., Muller, F., Lammers, F., Saleh, A., London, M., de Boer, P., \& Winterer, G. (2020). Ketamine effects on default mode network activity and vigilance: A randomized, placebo-controlled crossover simultaneous fMRI/EEG study. Human Brain Mapping, 41(1), 107-119. https://doi.org/https://doi.org/10.1002/hbm.24791 ALIUS Bulletin n6 (2024) aliusresearch.org/bulletin\par
  \endgroup
}%
\end{document}
